\CatchFileBetweenTags{\SpectralIndexVal}{\PMcalculations}{SpectralIndexVal}
\CatchFileBetweenTags{\SpectralIndexExperimentalValue}{\PMcalculations}{SpectralIndexExperimentalValue}
\CatchFileBetweenTags{\SpectralIndexAccText}{\PMcalculations}{SpectralIndexAccText}
\CatchFileBetweenTags{\JMatrixVal}{\PMcalculations}{JMatrixVal}
\CatchFileBetweenTags{\JGeometricVal}{\PMcalculations}{JGeometricVal}
\CatchFileBetweenTags{\EtaBaryonVal}{\PMcalculations}{EtaBaryonVal}
\CatchFileBetweenTags{\EtaBaryonExperimentalValue}{\PMcalculations}{EtaBaryonExperimentalValue}
\CatchFileBetweenTags{\EtaBaryonAccText}{\PMcalculations}{EtaBaryonAccText}
\CatchFileBetweenTags{\EtaBaryonGeoVal}{\PMcalculations}{EtaBaryonGeoVal}
\CatchFileBetweenTags{\EtaBaryonGeoSigma}{\PMcalculations}{EtaBaryonGeoSigma}
\CatchFileBetweenTags{\zeqVal}{\PMcalculations}{zeqVal}
\CatchFileBetweenTags{\zeqExperimentalValue}{\PMcalculations}{zeqExperimentalValue}
\CatchFileBetweenTags{\zeqAccText}{\PMcalculations}{zeqAccText}

\section{Primordial Initialization: Candidate Rules and Missing Dynamics}
\subsection{Chiral selection and reversible bookkeeping}
The channel count $N=2\nu=32$ and selected chiral count $\nu=16$ motivate a proposed initialization transition. These counts alone do not derive inflation, accelerated expansion or an arrow of time. There is also a reversibility constraint: if 32 distinct labels must remain distinguishable, they cannot be injectively mapped into only 16 labels. Retaining a binary record is sufficient, since a 32-label set can be relabeled as 16 labels times two record values. Lean proves both statements. Channel dimensions are not thereby identified with the number of complete microscopic states; the physical state and record spaces must be specified.

\subsection{Spectral-index ansatz}
The channel-overhead term suggests the candidate
\begin{equation}
n_s^{\rm ansatz}=1-S_Q=1-\frac{1}{2\nu}=\frac{31}{32}=\SpectralIndexVal.
\end{equation}
It is the complement of the same channel term used in the load budget. Compared with the draft reference \SpectralIndexExperimentalValue, it \SpectralIndexAccText. This arithmetic does not compute a primordial two-point function or its logarithmic slope. A perturbation-generation mechanism and a scale-dependent spectrum are needed to identify this quantity with the observed tilt.

\subsection{Tensor modes and initial conditions}
Setting the primordial tensor-to-scalar ratio to $r=0$ solely because the initialization phase has no matter load would require an additional argument. A vanishing source is insufficient to establish this conclusion: homogeneous modes depend on initial conditions or a fluctuation state. Lean illustrates the logical issue with $h(t)=\cos t$, which solves the source-free equation $h''+h=0$ and has $h(0)=1$. This is a mathematical counterexample to the inference from zero forcing to zero field, not a derived cosmological tensor equation. The substrate tensor dynamics and primordial state remain to be supplied before a prediction for $r$ can be made.

\subsection{Baryogenesis: distinguish the CP inputs}
The proposed conversion and dilution law is
\begin{equation}
\begin{split}
\eta_B(J)&=\frac{J}{\nu}
\frac{1}{(\pi\Delta)(H_{\rm sys}+\chi/D)}\\
&=\frac{J}{16168\pi},\qquad H_{\rm sys}=23.
\end{split}
\end{equation}
This is an operational hypothesis, not a consequence of a nonzero CP invariant. A baryon-number-changing transition, an initialization state and the subsequent dilution history must justify its factors. A conserved baryon label stays conserved under every iterate even if a separate CP parameter is nonzero; Lean checks that boundary explicitly.

Paper IV's complete unitary CKM ansatz has invariant $J_{\rm CKM}\approx\JMatrixVal$. The older independent geometric proposal gives $J_{\rm geo}\approx\JGeometricVal$. The two quantities are unequal in the current formalization and must not be substituted for each other. Applying the same conversion law gives
\begin{align}
\eta_B(J_{\rm CKM})&\approx\EtaBaryonVal,\\
\eta_B(J_{\rm geo})&\approx\EtaBaryonGeoVal.
\end{align}
For the draft reference \EtaBaryonExperimentalValue, the matrix-based value \EtaBaryonAccText; the geometric value differs by approximately $\EtaBaryonGeoSigma\sigma$. These are fixed-input residuals under the candidate law, not a combined statistical test or a reason to choose the closer CP input. The matrix value here is the full unitary invariant before any visible-sector dressing; such a readout correction would be an additional physical choice. The angle formulas and CP-to-baryon conversion need independent justification.

\subsection{Matter--radiation equality}
The geometric candidate is
\begin{equation}
z_{\rm eq}^{\rm ansatz}=(\pi\Delta)(H_{\rm sys}+\chi)
=1075\pi\approx\zeqVal.
\end{equation}
Against the draft reference \zeqExperimentalValue, it \zeqAccText. The geometric product alone does not establish the equality epoch. Under the separately supplied dilution laws
\begin{equation}
\rho_m(z)=\rho_{m0}(1+z)^3,\qquad
\rho_r(z)=\rho_{r0}(1+z)^4,
\end{equation}
positive $\rho_{r0}$ and $z>-1$ give
\begin{equation}
1+z_{\rm eq}=\frac{\rho_{m0}}{\rho_{r0}}.
\end{equation}
Lean proves this equivalence. Thus the proposed $z_{\rm eq}=1075\pi$ requires $\rho_{m0}/\rho_{r0}=1+1075\pi$. Neither that amplitude ratio nor the dilution dynamics follows from the slot product alone.
